\subsection{The actual three-plus-three iteration fails on the whole tested tensor}

The positive containing tensor in \fid{F66} does not
extend to every actual iterated tensor product.
Here the strand rank remains \(r=3\), while the
iteration parameter counts ordinary Hecke tensor
factors. Both input modules below have effective
two-state coefficient kernels. The obstruction
occurs on their entire four-dimensional tensor
product, without a projector inside that output.

Write
\[
 \delta=v+v^{-1},\quad D=\delta^2,\quad
 m=\delta^3,\quad a=3D-4,\quad d=m^2=D^3.
\]
The source variable \(u\) is renamed \(v\).
For the actual iterated algebra, Definition~2.5
and equations~(8)--(11), p.~6 of
\cite{BlasiakChebyshev} give
\begin{equation}\label{eq:sixTensorActualRestriction}
 \Delta^{(k_l,k_r)}(P_i^{(k_l+k_r)})
   =P_i^{(k_l)}\otimes P_i^{(k_r)}
       +Q_i^{(k_l)}\otimes Q_i^{(k_r)},
 \qquad Q_i^{(k)}=\delta^k-P_i^{(k)}.
\end{equation}
The following discussion on p.~7 supplies the
restricted module for each pair of modules.
This is the published algebra inclusion; no
Hopf coproduct on the whole ordinary Hecke
algebra is assumed.

\begin{finding}[\fid{F67}, an ineffective whole iterated tensor]
\label{f:sixTensorWholeObstruction}
The actual simple module \(\check M^{(3)}\) has
a generic basis in which its two named generators
have entries in \(\N[v,v^{-1}]\), and effective
graded matrix kernels with the correct quadratic
coefficient \(\delta^3\).

Nevertheless, the actual four-dimensional module
\(\check M^{(3)}\otimes\check M^{(3)}\), restricted
through \(\Delta^{(3,3)}\), has no basis making
both named generators nonnegative integral
Laurent matrices. Its unmarked two-word trace
has coefficient \(-12\) at each of \(v^{10}\)
and \(v^{-10}\).

Consequently no effective finite graded-kernel
iteration containing this input and preserving
the specified actual output, compositions and
unmarked word trace can realize this restriction.
Under either \(v=Q^{n/2}\) or \(v=Q^n\), the
same trace is not the Euler trace of any bounded
pure point complex of one integral weight.
A bounded mixed Tate complex does realize the
scalar trace.
\end{finding}

\begin{proof}
\par\smallskip\noindent\textbf{The actual effective inputs.}\quad
Equation~(13), p.~8 of \cite{BlasiakChebyshev}
defines \(\check X^{(k)}(c)\) in the left-column
convention by
\[
 P_1^{(k)}\longmapsto
   \begin{pmatrix}\delta^k&c\\0&0\end{pmatrix},
 \qquad
 P_2^{(k)}\longmapsto
   \begin{pmatrix}0&0\\1&\delta^k\end{pmatrix}.
\]
Theorem~3.4, p.~12 identifies
\[
 \check M^{(3)}
 =\check X^{(3)}\bigl((\delta^3T_3(1/\delta))^2\bigr),
\]
and equation~(24), p.~11 gives
\(\delta^3T_3(1/\delta)=4-3D\).
Thus its parameter is \(c=a^2\).
If \(S=\operatorname{diag}(a,1)\), replacing
each source matrix \(A_i\) by \(S^{-1}A_iS\)
gives the actual named matrices
\begin{equation}\label{eq:sixTensorPositiveInputs}
 N_1=\begin{pmatrix}m&a\\0&0\end{pmatrix},\qquad
 N_2=\begin{pmatrix}0&0\\a&m\end{pmatrix}.
\end{equation}
The source parameter uses a square, so the
choice \(a=3D-4\), opposite in sign to the
printed Chebyshev value, is legitimate.
The irreducibility criterion (14), p.~9 applies:
\(a^2\ne0,m^2\) over \(K=\Q(v)\), since
\begin{equation}\label{eq:sixTensorSimpleFactor}
 m^2-a^2=D^3-(3D-4)^2=(D-1)(D-4)^2\ne0.
\end{equation}

Both relevant characters are effective:
\[
 m=v^3+3v+3v^{-1}+v^{-3},\qquad
 a=3v^2+2+3v^{-2}.
\]
Choose finite graded spaces \(M,A\) with these
characters and form
\[
 \mathcal N_1=\begin{pmatrix}M&A\\0&0\end{pmatrix},\qquad
 \mathcal N_2=\begin{pmatrix}0&0\\A&M\end{pmatrix}.
\]
Matrix convolution uses direct sums and coefficient
tensor products. Its entries give actual
distributivity and symmetry isomorphisms
\begin{equation}\label{eq:sixTensorInputQuadratic}
              \mathcal N_i^2\simeq M\otimes\mathcal N_i.
\end{equation}
Thus these are effective input models with
the correct named action and quadratic coefficient.
No additional source-prescribed higher categorical
relations or canonical simple labels are inferred.

They also admit the full-Tate realization of
\fid{F61}: for a prime power \(Q\),
\(\ell\nmid Q\), and finite \(E/\Q_\ell\),
send a grade-\(e\) line to \(E(-e)[-2e]\)
on \(\Spec\F_Q\).
All shifts are even, so this is symmetric
monoidal and pure of complex weight zero,
with trace parameter \(v=Q^n\).
It applies entrywise to the displayed kernels
and their actual quadratic maps.
This uses ordinary Tate lines, and does not
assert a common spectral radius on total cohomology.

\par\smallskip\noindent\textbf{The full output matrices.}\quad
Apply \eqref{eq:sixTensorActualRestriction}
with \(k_l=k_r=3\), using the tensor basis
of \eqref{eq:sixTensorPositiveInputs}.
The entire four-dimensional named action is
\begin{equation}\label{eq:sixTensorActualMatrices}
 R_1=\begin{pmatrix}
 d&ma&ma&2a^2\\0&0&0&-ma\\0&0&0&-ma\\0&0&0&d
 \end{pmatrix},\qquad
 R_2=\begin{pmatrix}
 d&0&0&0\\-ma&0&0&0\\-ma&0&0&0\\2a^2&ma&ma&d
 \end{pmatrix}.
\end{equation}
Indeed
\[
 R_i=N_i\otimes N_i+(mI-N_i)\otimes(mI-N_i),
 \qquad R_i^2=dR_i.
\]
All four diagonal entries of \(R_1R_2\) are
\begin{equation}\label{eq:sixTensorFullDiagonal}
 m^4-2m^2a^2+4a^4,\quad -m^2a^2,\quad
              -m^2a^2,\quad m^4.
\end{equation}
Their sum is
\begin{equation}\label{eq:sixTensorTraceIdentity}
 \tau_6(v)=\Tr(R_1R_2)=d^2+A_6^2,
 \qquad
 A_6=2a^2-d=-D^3+18D^2-48D+32.
\end{equation}
Expanding \(D=v^2+2+v^{-2}\) gives
\begin{equation}\label{eq:sixTensorNegativeTrace}
\begin{split}
 \tau_6(v)={}&2(v^{12}+v^{-12})-12(v^{10}+v^{-10})
                  +192(v^8+v^{-8})\\
 &+388(v^6+v^{-6})+1134(v^4+v^{-4})
                  +1416(v^2+v^{-2})+1952.
\end{split}
\end{equation}
This exact product calculation proves both
negative coefficients without any simple projection.
The double-angle identity also gives
\(A_6=\delta^6T_6(1/\delta)\).
Proposition~3.1, equations~(16)--(17), p.~9
of \cite{BlasiakChebyshev} identifies the generic
decomposition as
\[
 \check\epsilon_+\oplus\check\epsilon_-
                 \oplus\check X^{(6)}(A_6^2).
\]
That decomposition is not needed for the
full-product trace proof.

\par\smallskip\noindent\textbf{The effective-operation obstruction.}\quad
Matrices over \(\N[v,v^{-1}]\) remain in that
semiring under multiplication, and their ordinary
trace has nonnegative coefficients.
Trace is invariant under every simultaneous
generic basis change, contradicting
\eqref{eq:sixTensorNegativeTrace}.
For finite graded kernels, the same argument
uses actual tensor products and finite direct
sums: the closed-word trace object would have
an effective graded character. It cannot have
the displayed character.

It follows that an effective iteration containing
the input \(\check M^{(3)}\), realizing the actual
restriction on its tensor square, and preserving
the named compositions and unmarked trace fails
on this example. Effective inputs alone do not
supply such an iteration. The source formula
also uses the complementary generators \(mI-N_i\),
which are virtual in the effective character cone.
The two negative terms in
\eqref{eq:sixTensorFullDiagonal} do not cancel
to an effective output trace.

Unlike \fid{F63}, no effective simple projector
inside the tested output is required.
Unlike the positive \(k=2+2\) model of \fid{F66},
the whole \(k=3+3\) tensor already fails.
This is an obstruction in the iterated tensor-factor
parameter, not an all-strand-rank binary theorem.
A category omitting this source simple as an
admissible input, or not comparing this actual
output, is not excluded.

\par\smallskip\noindent\textbf{Generic occurrence and the specialization limit.}\quad
The proof of Theorem~3.4, equations~(25)--(26),
p.~12 of \cite{BlasiakChebyshev}, realizes
\(\check M^{(2)}\) as a constituent of
\(\check M^{(1)}\otimes\check M^{(1)}\), and
\(\check M^{(3)}\) as a constituent of
\(\check M^{(1)}\otimes\check M^{(2)}\).
There \(\check M^{(1)}\) is the ordinary
Specht module of shape \((2,1)\).
Generic split semisimplicity in Corollary~3.5,
p.~13 therefore makes the tested tensor an
actual generic direct summand of the sixfold
ordinary Specht tensor \(M_{(2,1)}^{\otimes6}\).
The degree-three ordinary Schur--Weyl decomposition,
used explicitly in \fid{F21}/\fid{F43}, then
places it generically inside
\[
 (V_1\otimes\cdots\otimes V_6)^{\otimes3},
          \qquad \dim V_i=2,
\]
whose dimension is \(2^{18}\).
We have computed the trace of the selected
four-dimensional tensor, not of that whole
ambient module.

The chosen input matrices are regular at \(v=1\),
but their lattice cannot be identified with an
ambient direct-summand lattice at that point.
Indeed \(m=a=8\) there, and
\[
 G_1=2N_1/m-I=\begin{pmatrix}1&2\\0&-1\end{pmatrix},
 \qquad
 G_2=2N_2/m-I=\begin{pmatrix}-1&0\\2&1\end{pmatrix}
\]
satisfy \(G_i^2=I\), whereas
\begin{equation}\label{eq:sixTensorSpecializationBoundary}
             (G_1G_2)^3
               =\begin{pmatrix}7&6\\-6&-5\end{pmatrix}\ne I.
\end{equation}
In the ordinary ambient specialization, source
equation~(8) gives
\(P_i^{(3)}=4(1+s_i\otimes s_i\otimes s_i)\).
Its corresponding \(G_i\) are diagonal \(S_3\)
generators, whose product has cube identity.
Thus this chosen two-state specialization does
not factor through that ambient group image;
no regular basis change can identify it with
an ambient specialized direct summand.
The generic occurrence and Laurent-trace
obstruction do not require such an identification.

\par\smallskip\noindent\textbf{Pure and mixed point traces.}\quad
Under \(v=Q^{n/2}\), all exponents in
\eqref{eq:sixTensorNegativeTrace} are even, so
the distinct visible roots have even point weights.
Under \(v=Q^n\), their point weights are again
even. For a pure weight-\(w\) point complex,
BBD~5.1.8, pp.~126--127 permits a root of
point weight \(e\) only in degree \(e-w\)
\cite{BBD}. Hence all visible Euler multiplicities
here must have the same sign \((-1)^w\).
Eventual finite power-sum uniqueness contradicts
the mixed signs in
\eqref{eq:sixTensorNegativeTrace}.
Jordan blocks do not alter the multiplicities.
Hidden roots cannot cancel in different degrees:
each fixed root has only one permitted degree
in a pure complex.

Even without purity, the same power-sum uniqueness
excludes a genuine finite uninserted Weil trace,
or a single overall cohomological parity.
An overall Tate or parity shift does not
remove mixed coefficient signs.
Inserted-operator traces and bounded mixed
complexes are outside these exclusions.

For an explicit mixed witness under
\(v=Q^{n/2}\), let \(E_+\) contain the
positive-coefficient lines \(E(-e/2)\) of
\eqref{eq:sixTensorNegativeTrace}, with their
displayed multiplicities, in degree zero.
It has dimension \(8216\).
Put
\[
 E_-=E(-5)^{\oplus12}\oplus E(5)^{\oplus12}.
\]
Then \(E_+\oplus E_-[-1]\) has the required
Euler trace at every extension, with Euler
dimension \(8216-24=8192=\tau_6(1)\).
It is a scalar mixed Tate witness, not a
positive kernel realization or an
operation-compatible replacement for the
actual tensor action.
\end{proof}

Although its Laurent coefficients have both
signs, \(\tau_6(v)=D^6+A_6^2\) is positive
for positive real \(v\). The obstruction is
to effective Laurent characters and the
stated pure-point comparison, not a claim
of negative numerical traces there.

% Separate paragraph for Part III.
\paragraph{The \fid{F67} whole iterated tensor in Grothendieck's language.}
The actual named action and restriction are
finite-dimensional module data over \(\Q(v)\).
EGA~I, Theorem~1.3.7, p.~86,
Corollary~1.3.8, p.~87 and
Theorem~1.4.1, pp.~90--92 supply the affine
module dictionary \cite{EGA}.
SGA~3, I, 4.7.2--3 and cote~90, scan~17
supply action language \cite{SGA3,Archive};
they do not turn this three-strand algebra
into a Hopf algebra or supply an effective
categorical iteration.

Finite graded coefficient spaces have effective
split-torus characters in \(\N[v,v^{-1}]\).
SGA~6, V, \S\S2--3 and cote~133, scan~18
supply the effective/virtual representation
language \cite{SGA6,Archive}, with no
\(\lambda\)-ring assertion for the
noncommutative generator algebra.
The effective input kernels and the ineffective
actual output trace are distinguished by the
explicit calculation above.

The bounded point complexes fit SGA~6, I,
Definition~4.7, pp.~121--122 and
Example~4.8, p.~122 \cite{SGA6}.
BBD~5.1.5 and~5.1.8, pp.~125--127 fix the
integral-weight and pure-complex convention
\cite{BBD}.
The additional pure-word comparison is not
a published VIII axiom. The explicit mixed
Tate witness prevents an inference that
the scalar coefficient has no geometry.
